\documentclass[10pt]{article}

% Compilar con LuaLaTeX o XeLaTeX
\usepackage{fontspec}
\setmainfont{TeX Gyre Pagella} % si no está, cámbiala
\setsansfont{Fira Sans}       % instalar o cambiar
\setmonofont{DejaVu Sans Mono}       % instalar o cambiar

\usepackage[spanish]{babel}
\usepackage{microtype}
\usepackage{geometry}
\geometry{margin=2.5cm}
\usepackage{xcolor}
\definecolor{unabblue}{HTML}{003366}
\definecolor{unabgold}{HTML}{C69214}
\definecolor{lightblue}{HTML}{E6F0FF}

\usepackage[most]{tcolorbox}
\tcbset{enhanced,
  colback=lightblue!30,
  colframe=unabblue,
  fonttitle=\sffamily\bfseries\large,
  boxrule=1pt,arc=3mm,
  left=2mm,right=2mm,top=1mm,bottom=1mm,
  before skip=8pt,after skip=12pt
}

\usepackage{fancyhdr}
\pagestyle{fancy}
\fancyhf{}
\lhead{\sffamily\bfseries\textcolor{unabblue}{Prueba de Programación -- Versión A}}
\rhead{\sffamily\textcolor{unabblue}{PCFI161}}
\cfoot{\thepage}

\usepackage{amsmath,amssymb}
\usepackage{hyperref}
\hypersetup{colorlinks=true,linkcolor=unabblue,urlcolor=unabblue}


% ------------------------------------------------------------------------------- 
% ---------------------------------- Documento ----------------------------------
% ------------------------------------------------------------------------------- 

\begin{document}

\begin{center}
  {\huge\bfseries\textcolor{unabblue}{Solemne 01 de Programación}}\\[4pt]
  {\Large\sffamily\textcolor{unabgold}{Depto de Física y Astronomía}}\\[6pt]
  {\large\sffamily Versión A}\\[10pt]
  \today
\end{center}

\vspace{1em}
\noindent\textbf{Instrucciones}: Responda cada pregunta escribiendo un programa en \texttt{Python 3}. Use solo módulos de la biblioteca estándar. Comente brevemente su código.

\noindent\makebox[\linewidth]{\rule{\textwidth}{0.4pt}}

% ------------------------------------------------------------------------------- 
% --------------------------------- Pregunta 1 ----------------------------------
% ------------------------------------------------------------------------------- 

\begin{tcolorbox}[title={Pregunta 1 – Aproximación de $e^x$}]
La función exponencial $e^x$ se puede aproximar mediante la serie de Taylor:
\[
  e^x = 1 + x + \frac{x^2}{2!} + \frac{x^3}{3!} + \dots
\]
Escriba un programa que:
\begin{itemize}
  \item Pida al usuario un valor real $x$.
  \item Pida un entero positivo $n$ (número de términos).
  \item Calcule la suma de los primeros $n$ términos usando un ciclo \texttt{for}.
  \item Compare el resultado con \texttt{math.exp(x)}.
\end{itemize}
\end{tcolorbox}

% ------------------------------------------------------------------------------- 
% --------------------------------- Pregunta 2 ----------------------------------
% ------------------------------------------------------------------------------- 

\begin{tcolorbox}[title={Pregunta 2 – Ecuación cúbica}]
Considere la ecuación $ax^{3}+bx^{2}+cx+d=0$. 
Escriba un programa que:
\begin{itemize}
  \item Pida los coeficientes $a,\,b,\,c$.
  \item Genere el coeficiente entero $d$ al azar entre $-20$ y $20$.
  \item Implemente una función \texttt{resolver\_cubica(a,b,c,d)} que calcule las raíces reales usando la librería \texttt{numpy.roots}.
  \item Clasifique el número de raíces reales usando un \texttt{if}.
\end{itemize}
\end{tcolorbox}


% ------------------------------------------------------------------------------- 
% --------------------------------- Pregunta 3 ----------------------------------
% ------------------------------------------------------------------------------- 

\begin{tcolorbox}[title={Pregunta 3 – Clasificación de estrellas por luminosidad}]
La luminosidad relativa de una estrella de masa $M$ (en masas solares) se estima por
\[
  L = M^{3.5}
\]
Escriba un programa que:
\begin{itemize}
  \item Pida la masa $M$ de una estrella.
  \item Implemente una función \texttt{luminosidad(M)} que calcule $L$.
  \item Use un ciclo \texttt{while} para procesar varios valores hasta que el usuario ingrese \texttt{-1}.
  \item Use \texttt{if/else} para indicar si la estrella es más o menos luminosa que el Sol. 
  
\end{itemize}
\end{tcolorbox}

\end{document}

